\documentclass{beamer}
\usecolortheme{dove}
\setbeamertemplate{navigation symbols}{}
\usepackage{amsmath,amssymb,amsfonts,amsthm, multicol, subfigure, color}
\usepackage{bm}
\usepackage{graphicx}
\usepackage{tabularx}
\usepackage{booktabs}
\usepackage{hyperref}
\usepackage{pdfpages}
\usepackage{xcolor}
\definecolor{seagreen}{RGB}{46, 139, 87}
\definecolor{ucla}{RGB}{39, 116, 174}
\definecolor{darkestblue}{RGB}{0, 59, 92}
\definecolor{gold}{RGB}{255, 209, 0}
\def\independenT#1#2{\mathrel{\rlap{$#1#2$}\mkern2mu{#1#2}}}
\newcommand\indep{\protect\mathpalette{\protect\independenT}{\perp}}
\def\log{\text{log}}
\newcommand\logit{\text{logit}}
\newcommand\iid{\stackrel{\text{iid}}{\sim}}
\newcommand\E{\text{E}}
\newcommand\V{\text{V}}
\renewcommand\P{\text{P}}
\newcommand{\Cov}{\text{Cov}}
\newcommand{\Cor}{\text{Cor}}
\newcommand\doop{\text{do}}


\usepackage{stackrel}
\usepackage{tikz}
\usetikzlibrary{arrows,shapes.arrows,positioning,shapes,patterns,calc}
\newcommand\slideref[1]{\vskip .1cm \tiny \textcolor{gray}{{#1}}}
\newcommand\red[1]{\color{red}#1}
\newcommand\blue[1]{\color{blue}#1}
\newcommand\gray[1]{\color{gray}#1}
\newcommand\seagreen[1]{\color{seagreen}#1}
\newcommand\purple[1]{\color{purple}#1}
\newcommand\orange[1]{\color{orange}#1}
\newcommand\black[1]{\color{black}#1}
\newcommand\white[1]{\color{white}#1}
\newcommand\teal[1]{\color{teal}#1}
\newcommand\magenta[1]{\color{magenta}#1}
\newcommand\Fuchsia[1]{\color{Fuchsia}#1}
\newcommand\BlueGreen[1]{\color{BlueGreen}#1}
\newcommand\bblue[1]{\textcolor{blue}{\textbf{#1}}}
\newcommand\bred[1]{\textcolor{red}{\textbf{#1}}}
\newcommand\bgray[1]{\textcolor{gray}{\textbf{#1}}}
\newcommand\bgreen[1]{\textcolor{seagreen}{\textbf{#1}}}
\newcommand\bref[2]{\href{#1}{\color{blue}{#2}}}
\colorlet{lightgray}{gray!40}
\pgfdeclarelayer{bg}    % declare background layer for tikz
\pgfsetlayers{bg,main} % order layers for tikz
\newcommand\mycite[1]{\begin{scriptsize}\textcolor{darkgray}{(#1)}\end{scriptsize}}
\newcommand{\tcframe}{\frame{
%\small{
\only<1|handout:0>{\tableofcontents}
\only<2|handout:1>{\tableofcontents[currentsection]}}
%}
}

\setbeamertemplate{footline}[frame number]


\usepackage[round]{natbib}
\bibliographystyle{humannat-mod}
\setbeamertemplate{enumerate items}[default]
\usepackage{mathtools}

\newcommand{\goalsframe}{\begin{frame}{Learning goals for today}
At the end of class, you will be able to:
\begin{enumerate}
\item Explain exchangeability in randomized experiments
\item Make arguments about why exchangeability may not hold in particular cases when a treatment is not randomized
\end{enumerate} \vskip .2in
\end{frame}}

\title{Randomized Experiments\footnote{Some material in this lecture draws on past materials by Sam Wang at Cornell University. Thanks Sam!}}
\author{Sociol 114}
\date{30 Jan 2025}

\begin{document}

\maketitle

\goalsframe

\begin{frame}{Exchangeable sampling from a population}
\pause
\begin{tikzpicture}[x = \textwidth, y = .9\textheight]
% POPULATION
\node[anchor = north, font = \bf, color = ucla, align = center] at (.15,1) {Population\\Outcomes};
\foreach \i in {.3,.4,.5,.6,.7,.8} {
	\draw[fill = ucla, draw = white] (.08,\i - .05) rectangle (.22,\i + .05) {};
}
\node[font = {\bf}, white] at (.15,.8) {$Y_\text{Maria}$};
\node[font = {\bf}, white] at (.15,.7) {$Y_\text{William}$};
\node[font = {\bf}, white] at (.15,.6) {$Y_\text{Rich}$};
\node[font = {\bf}, white] at (.15,.5) {$Y_\text{Sarah}$};
\node[font = {\bf}, white] at (.15,.4) {$Y_\text{Alondra}$};
\node[font = {\bf}, white] at (.15,.3) {$Y_\text{Jes\'us}$};
\pause
% RANDOMIZATION
\node[anchor = north, font = \bf, color = ucla, align = center] at (.4,1) {Randomized\\Sampling};
%\foreach \i in {.3,.4,.5,.6,.7,.8} {
%	\draw[fill = ucla, draw = white] (.4,\i) rectangle (6,\i + .1) {};
%}
\node[font = {\bf}, anchor = east, ucla] at (.5,.8) {$S_\text{Maria} = 1$};
\node[font = {\bf}, anchor = east, ucla] at (.5,.7) {$S_\text{William} = 0$};
\node[font = {\bf}, anchor = east, ucla] at (.5,.6) {$S_\text{Rich} = 0$};
\node[font = {\bf}, anchor = east, ucla] at (.5,.5) {$S_\text{Sarah} = 1$};
\node[font = {\bf}, anchor = east, ucla] at (.5,.4) {$S_\text{Alondra} = 0$};
\node[font = {\bf}, anchor = east, ucla] at (.5,.3) {$S_\text{Jes\'us} = 1$};
\pause
% SAMPLE
\node[anchor = north, font = \bf, color = ucla, align = center] at (.62,1) {Sampled\\Outcomes};
\foreach \i in {.3,.5,.8} {
	\draw[fill = ucla, draw = white] (.55,\i - .05) rectangle (.69,\i + .05) {};
}
\node[font = {\bf}, white] at (.62,.8) {$Y_\text{Maria}$};
%\node[font = {\bf}, white] at (.75,.75) {$Y_\text{William}$};
%\node[font = {\bf}, white] at (.75,.65) {$Y_\text{Rich}$};
\node[font = {\bf}, white] at (.62,.5) {$Y_\text{Sarah}$};
%\node[font = {\bf}, white] at (.75,.45) {$Y_\text{Alondra}$};
\node[font = {\bf}, white] at (.62,.3) {$Y_\text{Jes\'us}$};
%\draw[fill = darkgray] (.75,0) rectangle (2,1);
\pause
\node[anchor = north west, align = left] at (.75,1) {\textbf{Estimator:}\\Estimate the\\population mean\\by the sample mean};
\node[anchor = north west, align = left] at (.75,.65) {\textbf{Key assumption}:\\Sampled and\\unsampled units\\are \textbf{exchangeable}\\due to random\\sampling};
\node[anchor = north west, align = left] at (.75,.27) {$Y\indep S$};
\end{tikzpicture}
\end{frame}

\begin{frame}
Now suppose our population all participate\\in a randomized experiment with\\treatment ($A = 1$) and control ($A = 0$) conditions
\end{frame}


\begin{frame}{Exchangeable treatment assignment} %Y1
\begin{tikzpicture}[x = \textwidth, y = .9\textheight]
% POPULATION
\node[anchor = north, font = \bf, color = ucla, align = center] at (.15,1.05) {Population\\Potential\\Outcomes};
\foreach \i in {.3,.4,.5,.6,.7,.8} {
	\draw[fill = ucla, draw = white] (.08,\i - .05) rectangle (.22,\i + .05) {};
}
\node[font = {\bf}, white] at (.15,.8) {$Y^1_\text{Maria}$};
\node[font = {\bf}, white] at (.15,.7) {$Y^1_\text{William}$};
\node[font = {\bf}, white] at (.15,.6) {$Y^1_\text{Rich}$};
\node[font = {\bf}, white] at (.15,.5) {$Y^1_\text{Sarah}$};
\node[font = {\bf}, white] at (.15,.4) {$Y^1_\text{Alondra}$};
\node[font = {\bf}, white] at (.15,.3) {$Y^1_\text{Jes\'us}$};
\pause
% RANDOMIZATION
\node[anchor = north, font = \bf, color = ucla, align = center] at (.4,1) {Randomized\\Treatment};
%\foreach \i in {.3,.4,.5,.6,.7,.8} {
%	\draw[fill = ucla, draw = white] (.4,\i) rectangle (6,\i + .1) {};
%}
\node[font = {\bf}, anchor = east, ucla] at (.5,.8) {$A_\text{Maria} = 1$};
\node[font = {\bf}, anchor = east, ucla] at (.5,.7) {$A_\text{William} = 0$};
\node[font = {\bf}, anchor = east, ucla] at (.5,.6) {$A_\text{Rich} = 0$};
\node[font = {\bf}, anchor = east, ucla] at (.5,.5) {$A_\text{Sarah} = 1$};
\node[font = {\bf}, anchor = east, ucla] at (.5,.4) {$A_\text{Alondra} = 0$};
\node[font = {\bf}, anchor = east, ucla] at (.5,.3) {$A_\text{Jes\'us} = 1$};
\pause
% SAMPLE
\node[anchor = north, font = \bf, color = ucla, align = center] at (.62,1) {Observed\\Outcomes};
\foreach \i in {.3,.5,.8} {
	\draw[fill = ucla, draw = white] (.55,\i - .05) rectangle (.69,\i + .05) {};
}
\node[font = {\bf}, white] at (.62,.8) {$Y^1_\text{Maria}$};
%\node[font = {\bf}, white] at (.75,.75) {$Y_\text{William}$};
%\node[font = {\bf}, white] at (.75,.65) {$Y_\text{Rich}$};
\node[font = {\bf}, white] at (.62,.5) {$Y^1_\text{Sarah}$};
%\node[font = {\bf}, white] at (.75,.45) {$Y_\text{Alondra}$};
\node[font = {\bf}, white] at (.62,.3) {$Y^1_\text{Jes\'us}$};
%\draw[fill = darkgray] (.75,0) rectangle (2,1);
\pause
\node[anchor = north west, align = left] at (.75,1) {\textbf{Estimator:}\\Estimate the\\population mean\\$\E(Y^1)$ by the\\untreated sample mean};
\node[anchor = north west, align = left] at (.75,.65) {\textbf{Key assumption}:\\Treated and\\untreated units\\are \textbf{exchangeable}\\due to random\\treatment assignment};
\node[anchor = north west, align = left] at (.75,.27) {$Y^1\indep A$};
\end{tikzpicture}
\end{frame}


\begin{frame}{Exchangeable treatment assignment} %Y0
\begin{tikzpicture}[x = \textwidth, y = .9\textheight]
% POPULATION
\node[anchor = north, font = \bf, color = ucla, align = center] at (.15,1.05) {Population\\Potential\\Outcomes};
\foreach \i in {.3,.4,.5,.6,.7,.8} {
	\draw[fill = gold, draw = white] (.08,\i - .05) rectangle (.22,\i + .05) {};
}
\node[font = {\bf}, darkestblue] at (.15,.8) {$Y^0_\text{Maria}$};
\node[font = {\bf}, darkestblue] at (.15,.7) {$Y^0_\text{William}$};
\node[font = {\bf}, darkestblue] at (.15,.6) {$Y^0_\text{Rich}$};
\node[font = {\bf}, darkestblue] at (.15,.5) {$Y^0_\text{Sarah}$};
\node[font = {\bf}, darkestblue] at (.15,.4) {$Y^0_\text{Alondra}$};
\node[font = {\bf}, darkestblue] at (.15,.3) {$Y^0_\text{Jes\'us}$};
% RANDOMIZATION
\node[anchor = north, font = \bf, color = ucla, align = center] at (.4,1) {Randomized\\Treatment};
%\foreach \i in {.3,.4,.5,.6,.7,.8} {
%	\draw[fill = ucla, draw = white] (.4,\i) rectangle (6,\i + .1) {};
%}
\node[font = {\bf}, anchor = east, ucla] at (.5,.8) {$A_\text{Maria} = 1$};
\node[font = {\bf}, anchor = east, ucla] at (.5,.7) {$A_\text{William} = 0$};
\node[font = {\bf}, anchor = east, ucla] at (.5,.6) {$A_\text{Rich} = 0$};
\node[font = {\bf}, anchor = east, ucla] at (.5,.5) {$A_\text{Sarah} = 1$};
\node[font = {\bf}, anchor = east, ucla] at (.5,.4) {$A_\text{Alondra} = 0$};
\node[font = {\bf}, anchor = east, ucla] at (.5,.3) {$A_\text{Jes\'us} = 1$};
% SAMPLE
\node[anchor = north, font = \bf, color = ucla, align = center] at (.62,1) {Observed\\Outcomes};
\foreach \i in {.4,.6,.7} {
	\draw[fill = gold, draw = white] (.55,\i - .05) rectangle (.69,\i + .05) {};
}
%\node[font = {\bf}, white] at (.62,.8) {$Y^1_\text{Maria}$};
\node[font = {\bf}, darkestblue] at (.62,.7) {$Y^0_\text{William}$};
\node[font = {\bf}, darkestblue] at (.62,.6) {$Y^0_\text{Rich}$};
%\node[font = {\bf}, white] at (.62,.5) {$Y^1_\text{Sarah}$};
\node[font = {\bf}, darkestblue] at (.62,.4) {$Y^0_\text{Alondra}$};
%\node[font = {\bf}, white] at (.62,.3) {$Y^1_\text{Jes\'us}$};
\node[anchor = north west, align = left] at (.75,1) {\textbf{Estimator:}\\Estimate the\\population mean\\$\E(Y^0)$ by the\\untreated sample mean};
\node[anchor = north west, align = left] at (.75,.65) {\textbf{Key assumption}:\\Treated and\\untreated units\\are \textbf{exchangeable}\\due to random\\treatment assignment};
\node[anchor = north west, align = left] at (.75,.27) {$Y^0\indep A$};
\end{tikzpicture}
\end{frame}

\begin{frame}{Exchangeable treatment assignment} %Y0 AND Y1
\begin{tikzpicture}[x = \textwidth, y = .9\textheight]
% POPULATION
\node[anchor = north, font = \bf, color = ucla, align = center] at (.23,1.05) {Population\\Potential\\Outcomes};
\foreach \i in {.3,.4,.5,.6,.7,.8} {
	\draw[fill = ucla, draw = white] (.08,\i - .05) rectangle (.22,\i + .05) {};
	\draw[fill = gold, draw = white] (.24,\i - .05) rectangle (.38,\i + .05) {};
}
\node[font = {\bf}, white] at (.15,.8) {$Y^1_\text{Maria}$};
\node[font = {\bf}, white] at (.15,.7) {$Y^1_\text{William}$};
\node[font = {\bf}, white] at (.15,.6) {$Y^1_\text{Rich}$};
\node[font = {\bf}, white] at (.15,.5) {$Y^1_\text{Sarah}$};
\node[font = {\bf}, white] at (.15,.4) {$Y^1_\text{Alondra}$};
\node[font = {\bf}, white] at (.15,.3) {$Y^1_\text{Jes\'us}$};
\node[font = {\bf}, darkestblue] at (.31,.8) {$Y^0_\text{Maria}$};
\node[font = {\bf}, darkestblue] at (.31,.7) {$Y^0_\text{William}$};
\node[font = {\bf}, darkestblue] at (.31,.6) {$Y^0_\text{Rich}$};
\node[font = {\bf}, darkestblue] at (.31,.5) {$Y^0_\text{Sarah}$};
\node[font = {\bf}, darkestblue] at (.31,.4) {$Y^0_\text{Alondra}$};
\node[font = {\bf}, darkestblue] at (.31,.3) {$Y^0_\text{Jes\'us}$};
% RANDOMIZATION
\node[anchor = north, font = \bf, color = ucla, align = center] at (.51,1) {Randomized\\Treatment};
%\foreach \i in {.3,.4,.5,.6,.7,.8} {
%	\draw[fill = ucla, draw = white] (.4,\i) rectangle (6,\i + .1) {};
%}
\node[font = {\bf}, anchor = east, ucla] at (.61,.8) {$A_\text{Maria} = 1$};
\node[font = {\bf}, anchor = east, ucla] at (.61,.7) {$A_\text{William} = 0$};
\node[font = {\bf}, anchor = east, ucla] at (.61,.6) {$A_\text{Rich} = 0$};
\node[font = {\bf}, anchor = east, ucla] at (.61,.5) {$A_\text{Sarah} = 1$};
\node[font = {\bf}, anchor = east, ucla] at (.61,.4) {$A_\text{Alondra} = 0$};
\node[font = {\bf}, anchor = east, ucla] at (.61,.3) {$A_\text{Jes\'us} = 1$};
% SAMPLE
\node[anchor = north, font = \bf, color = ucla, align = center] at (.83,1) {Observed\\Outcomes};
\foreach \i in {.3,.5,.8} {
	\draw[fill = ucla, draw = white] (.68,\i - .05) rectangle (.82,\i + .05) {};
}
\foreach \i in {.4,.6,.7} {
	\draw[fill = gold, draw = white] (.84,\i - .05) rectangle (.98,\i + .05) {};
}
\node[font = {\bf}, white] at (.75,.8) {$Y^1_\text{Maria}$};
\node[font = {\bf}, darkestblue] at (.91,.7) {$Y^0_\text{William}$};
\node[font = {\bf}, darkestblue] at (.91,.6) {$Y^0_\text{Rich}$};
\node[font = {\bf}, white] at (.75,.5) {$Y^1_\text{Sarah}$};
\node[font = {\bf}, darkestblue] at (.91,.4) {$Y^0_\text{Alondra}$};
\node[font = {\bf}, white] at (.75,.3) {$Y^1_\text{Jes\'us}$};
%\node[anchor = north west, align = left] at (.75,1) {\textbf{Estimator:}\\Estimate the\\population mean\\$\E(Y^0)$ by the\\untreated sample mean};
%\node[anchor = north west, align = left] at (.75,.65) {\textbf{Key assumption}:\\Treated and\\untreated units\\are \textbf{exchangeable}\\due to random\\treatment assignment};
%\node[anchor = north west, align = left] at (.75,.27) {$Y^0\indep A$};
\end{tikzpicture}
\end{frame}

\begin{frame}{Exchangeable treatment assignment}
\textbf{Causal Estimand:}\\(expected outcome if assigned to treatment) \\
${}-{}$ (expected outcome if assigned to control)
$$
\E\left(Y^1\right) - \E\left(Y^0\right)
$$ \vskip .1in
\textbf{Exchangeability Assumption}:\\
Potential outcomes are independent of treatment assignment \\
$$\{Y^0,Y^1\}\indep A$$ \vskip .1in
\textbf{Empirical Estimand:}\\
(expected outcome among the treated)\\
$-$ (expected outcome among the untreated)
$$\E(Y\mid A = 1) - \E(Y\mid A = 0)$$
\end{frame}

\begin{frame}{Exchangeable treatment assignment: Proof}
$$\begin{aligned}
&\E\left(Y^1\right) - \E\left(Y^0\right) \\
&= \E\left(Y^1\mid A = 1\right) - \E\left(Y^0\mid A = 0\right) &\onslide<3->{\text{by exchangeability}}\\
&= \E\left(Y\mid A = 1\right) - \E\left(Y\mid A = 0\right) &\onslide<2->{\text{by consistency}}
\end{aligned}$$ \vskip .2in \pause \pause \pause
This is an example of \textbf{causal identification}: \\
using assumptions to arrive at an empirical quantity\\(involving only factual random variables, no potential outcomes)\\that equals our causal estimand if the assumptions hold \vskip .2in
The causal estimand $\E(Y^1) - \E(Y^0)$ is \textbf{identified} by the empirical estimand $\E(Y\mid A = 1) - \E(Y\mid A = 0)$
\end{frame}

\begin{frame}{Potential outcome exercise: Covid vaccines} \pause
% from Sam
Suppose we know the following pieces of information:
{\small
\begin{itemize}
\item Martha was vaccinated against Covid.\\Martha tested negative 6 months later.
\item Ezra was not vaccinated against Covid.\\Ezra tested positive 6 months later.
\end{itemize} \pause
Which cells have known values? What are the values?
}

\begin{table}
  \renewcommand*{\arraystretch}{2}
\begin{tabular}[t]{c|c|c| c|c}
  & $A_i$ & $Y_i$ & $Y_i^{\text{Vaccinated}}$ & $Y_i^{\text{Unvaccinated}}$\\
\hline
Martha & \qquad \qquad \qquad & \qquad \qquad \qquad & \qquad & \qquad\\ \hline
Ezra & \qquad \qquad & \qquad \qquad & \qquad & \qquad\\
\end{tabular}
\end{table}

\end{frame}

\begin{frame}{Experiments vs observational studies}
Suppose we gathered data by surveying individuals in Fall of 2021
    \begin{itemize}
        \item Vaccinated for covid ($A_i = 1$) or not ($A_i = 0$)
        \item Tested positive for Covid in 2021: yes ($Y_i = 1$) or no ($Y_i=0$)
     \end{itemize}
\end{frame}

\begin{frame}{Experiments vs observational studies}
We observe evidence
     \begin{itemize}
         \item Of the individuals who are \textbf{vaccinated} ($A_i = 1$),\\50\% had a positive Covid test ($Y_i = 1$)
         \item Of the individuals who are \textbf{not vaccinated} ($A_i = 0$),\\70\% had a positive Covid test ($Y_i = 1$) 
\end{itemize} \pause
How might a vaccine skeptic explain the data? 
\end{frame}

\begin{frame}{Experiments vs observational studies}
Experiment designed by Pfizer \textbf{randomly assign} each individual (43,000 total) into two groups\footnote{Polack et. al. NEJM 2020}:
\begin{itemize}
        \item Two doses of BNT162b2 vaccine 21 days apart
        \item Two doses of placebo 21 days apart
        \item Whether a positive Covid test was recorded between 7 days and 14 weeks after the injection
\end{itemize}
\vspace{1em}

\pause
     \begin{itemize}
         \item Of the individuals who were given the vaccine ($A_i = 1$), 0.04\% had a positive Covid test ($Y_i = 1$)
         \item Of the individuals who were given the placebo ($A_i = 0$), 0.9\% had a positive Covid test ($Y_i = 1$)
         \item Individuals who received the placebo were $\approx 20$ times more likely to get Covid 
     \end{itemize}

\pause \vspace{1em}
\textbf{Do the skeptic's objections still hold?}

\end{frame}

\begin{frame}{Why experiments ``work'': Exchangeability}

\begin{figure}
    \centering
    \includegraphics[scale = .18]{figures/pfizer_balance.jpg}
\end{figure}

\end{frame}



\begin{frame}{Why experiments ``work'': Exchangeability}

In random experiments, the distribution of \textbf{potential outcomes} for those who are treated and those who are not treated (control group) are identically distributed!

$$\{Y^1,Y^0\}\indep A$$

\vspace{2em}
This is \textbf{exchangeability} \vskip .2in

\textbf{Question.} Does exchangeability imply $Y\indep A$?

\end{frame}

\begin{frame}{Why experiments ``work'': Exchangeability}

Exchangeability is about \textbf{potential} rather than \textbf{observed} outcomes
$$\{Y^0,Y^1\}\indep A\qquad\text{rather than}\qquad Y\not\indep A$$
 \pause
\begin{itemize}
\item Potential outcomes are independent of treatment $\{Y^0,Y^1\}\indep A$
\begin{itemize}
\item Example: Risk of covid under no vaccine ($Y^0$) is the same for those with and without a vaccine
\end{itemize} \pause
\item Observed outcomes are not independent of treatment $Y\not\indep A$
\begin{itemize}
\item Example: Risk of covid is lower for those with the vaccine
\item Why? Because for them $Y = Y^1$. For others, $Y = Y^0$.
\item If $A$ affects $Y$, then $Y\not\indep A$
\end{itemize}
\end{itemize} \pause
Under exchangeability, the only reason $Y\not\indep A$ is if $A$ causes $Y$.
\end{frame}

\begin{frame}{Design a hypothetical experiment}

\begin{itemize}
\item Define a treatment and an outcome
\item Design a randomized experiment
\begin{itemize}
\item Who would you enroll?
\item How would you randomize the treatment?
\item When and how would you measure the outcome?
\end{itemize}
\item Think of a criticism that could be levied against you if you had not randomized the treatment, which is overcome by randomization
\end{itemize}

\end{frame}

\goalsframe

\end{document}


\begin{frame}{Randomized experiments}
\begin{figure}
    \centering
    \includegraphics[scale = .3]{figures/goldStandard.png}
\end{figure}
    
\end{frame}

\begin{frame}{Fundamental problem of causal inference}
\begin{itemize}
     \item Randomized experiments are the gold standard for estimating causal effects
    \item Fundamental problem of causal inference is that we don't observe counterfactual outcomes
    \item Data is still missing in random experiments
\end{itemize}

\begin{table}\footnotesize
    \centering

\begin{tabular}[t]{c|c| c|c|c}
\toprule
  & A & $Y^{a=1}$ & $Y^{a=0}$ & $Y^{a=1} - Y^{a=0}$\\
\midrule
Ind 1 & 0 & ? & 0 & ?\\
Ind 2 & 0 & ? & 1 & ?\\
Ind 3 & 0  & ? & 0 & ?\\
Ind 4 &1 & 1 & ? & ?\\
Ind 5 & 1 & 0 & ? & ?\\
Ind 6 & 1 & 1 & ? & ?\\
\bottomrule
\end{tabular}
\end{table}

\pause 
\begin{itemize}
    \item Why do randomized experiments ``work''?
\end{itemize}
    
\end{frame}












\begin{frame}{\only<1-3>{The fundamental problem of causal inference}\only<4->{Preview: Solving the problem by assumptions}}
\begin{tikzpicture}[x = \textwidth, y = .8\textheight]
\node at (0,0) {};
\node at (1,1) {};
\node<1-3>[anchor = south east, gray, font = \bf] at (1,0) {\href{https://www.tandfonline.com/doi/abs/10.1080/01621459.1986.10478354}{Holland 1986}};
% Individual outcomes
\node[anchor = north west, font = \bf] at (.05,1) {The data};
\foreach \i in {.5,.7,.8} {
	\draw[fill = blue, opacity = .3, color = blue] (.05,\i) rectangle (.25,\i + .1) {};
}
\foreach \i in {.3,.4,.6} {
	\draw[fill = seagreen, opacity = .3, color = seagreen] (.25,\i) rectangle (.45,\i + .1) {};
}
\node[font = \footnotesize] at (.15,.85) {$Y_\text{Nick}^\text{College}$};
\node[font = \footnotesize] at (.15,.75) {$Y_\text{William}^\text{College}$};
\node[font = \footnotesize] at (.35,.65) {$Y_\text{Rich}^\text{Non-college}$};
\node[font = \footnotesize] at (.15,.55) {$Y_\text{Diego}^\text{College}$};
\node[font = \footnotesize] at (.35,.45) {$Y_\text{Javier}^\text{Non-college}$};
\node[font = \footnotesize] at (.35,.35) {$Y_\text{Jes\'us}^\text{Non-college}$};
\node[anchor = north, align = center, font = \footnotesize, blue] at (.15, .3) {Outcome\\under\\treatment};
\node[anchor = north, align = center, font = \footnotesize, seagreen] at (.35, .3) {Outcome\\under\\control};
\node[anchor = south, rotate = 90, align = center] at (.05, .6) {Each Row is a Person};
% The claim
\onslide<2->{
\node[anchor = north west, font = \bf] at (.55,1) {The claim};
\foreach \i in {.3,.4,.5,.6,.7,.8} {
	\draw[fill = blue, opacity = .3, color = blue] (.55,\i) rectangle (.75,\i + .1) {};
	\draw[fill = seagreen, opacity = .3, color = seagreen] (.75,\i) rectangle (.95,\i + .1) {};
	\draw[<->, thick] (.73,\i + .05) -- (.77,\i + .05);
}
\node[font = \footnotesize] at (.65,.85) {$Y_\text{Nick}^\text{College}$};
\node[font = \footnotesize] at (.85,.85) {$Y_\text{Nick}^\text{Non-college}$};
\node[font = \footnotesize] at (.65,.75) {$Y_\text{William}^\text{College}$};
\node[font = \footnotesize] at (.85,.75) {$Y_\text{William}^\text{Non-college}$};
\node[font = \footnotesize] at (.65,.65) {$Y_\text{Rich}^\text{College}$};
\node[font = \footnotesize] at (.85,.65) {$Y_\text{Rich}^\text{Non-college}$};
\node[font = \footnotesize] at (.65,.55) {$Y_\text{Diego}^\text{College}$};
\node[font = \footnotesize] at (.85,.55) {$Y_\text{Diego}^\text{Non-college}$};
\node[font = \footnotesize] at (.65,.45) {$Y_\text{Javier}^\text{College}$};
\node[font = \footnotesize] at (.85,.45) {$Y_\text{Javier}^\text{Non-college}$};
\node[font = \footnotesize] at (.65,.35) {$Y_\text{Jes\'us}^\text{College}$};
\node[font = \footnotesize] at (.85,.35) {$Y_\text{Jes\'us}^\text{Non-college}$};
\node[anchor = north, align = center, font = \footnotesize, blue] at (.65, .3) {Outcome\\under\\treatment};
\node[anchor = north, align = center, font = \footnotesize, seagreen] at (.85, .3) {Outcome\\under\\control};
}
\node<3>[anchor = north west] at (0,.06) {Counterfactuals are \textbf{not observed}};
\onslide<3->{
\node[font = \footnotesize] at (.35,.85) {?};
\node[font = \footnotesize] at (.35,.75) {?};
\node[font = \footnotesize] at (.15,.65) {?};
\node[font = \footnotesize] at (.35,.55) {?};
\node[font = \footnotesize] at (.15,.45) {?};
\node[font = \footnotesize] at (.15,.35) {?};
}
\draw<5-> (.05,.6) rectangle (.45,.9);
\draw<5-> (.05,.3) rectangle (.45,.6);
\draw<6->[->, thick] (.3, .68) to[out = 90, in = 215] (.33,.73);
\draw<6->[->, thick] (.3, .68) to[out = 90, in = 215] (.33,.83);
\draw<7->[->, thick] (.1, .83) to[out = 215, in = 150] (.13,.65);
\draw<7->[->, thick] (.1, .73) to[out = 215, in = 150] (.13,.65);
\draw<8->[->, thick] (.116, .53) to[out = 270, in = 150] (.14,.47);
\draw<8->[->, thick] (.116, .53) to[out = 270, in = 150] (.14,.37);
\draw<9->[->, thick] (.3, .47) to[out = 115, in = 210] (.33,.53);
\draw<9->[->, thick] (.28, .37) to[out = 145, in = 210] (.33,.53);
\end{tikzpicture}
\end{frame}

\section{Quick Refresher}

\begin{frame}{Potential outcome notation review}
\begin{itemize}
\item We typically use $i$ to denote a generic unit in our study
\item $Y_i$ is the observed outcome for unit $i$
\item $A_i$ is treatment received by unit $i$
\item We typically use $a$ to denote a generic treatment in our study
\item $Y^a_i$ denotes the outcome we would have observed for unit $i$ if assigned treatment value $a$
\end{itemize}
\end{frame}


\begin{frame}{Potential outcome exercise}
\begin{itemize}
\item We observe that Martha ate a Mediterranean diet, and we observe that Martha survived.
\\ Suppose Martha had eaten a standard diet, we would have observed that Martha survived. 
\item We observe that Ezra ate a standard diet, and we observe that Ezra did not survive
\\ Suppose Ezra had eaten a Mediterranean diet, we would have observed that Ezra survived. 
\end{itemize}
\pause

\vspace{1em}

Assuming \emph{consistency}, what is $A_i$, $Y_i$, $Y_i^{\texttt{MedDiet}}$ and $Y_i^{\texttt{StanDiet}}$
\begin{itemize}
\item When $i = \text{Martha}$?
\item When $i = \text{Ezra}$? 
\end{itemize}


\end{frame}





\section{Randomized experiments: Two key benefits}












\begin{frame}{Exchangeability}
\textbf{Definition:} Exchangeability means that the potential outcomes, $Y^{a=1}$ and $Y^{a = 0}$, are independent of the observed treatment ($A$)
\pause
\vspace{1em}

{\small 
\begin{itemize}
    \item Two random variables are \textbf{independent} if the outcome of one does not give any information about the outcome of the other
    \pause 
    \item What we would have observed if an individual was given the treatment ($Y^{a = 1}$) is independent of whether or not the individual actually received treatment 
    \item What we would have observed if an individual was not given the treatment ($Y^{a = 0}$) is independent of whether or not the individual actually received treatment  
    \pause 
     \item Exchangeability means that the \textbf{potential} outcomes $Y^a$ are independent of the observed treatment
    \item Exchangeability does \textbf{not} mean that the \textbf{observed} outcome $Y$ is independent of the observed treatment!
\end{itemize}
}


\end{frame}


\begin{frame}{Exchangeability}
$A = 1$ means vaccinated; $A = 0$ means unvaccinated\\
$Y = 1$ means covid; $Y = 0$ means no covid;

\begin{table}
\begin{tabular}[t]{c|c| c|c|c}
\toprule
& $Y^{a=1}$ & $Y^{a=0}$ & $A$ & $Y$\\
\midrule
Low Risk 1 & 0 & 0 & ? & ?\\
Low Risk 2 & 0 & 0 & ? & ?\\
High Risk 3 & 0 & 1 & ? & ?\\
High Risk 4 & 0 & 1 & ? & ?\\
\bottomrule
\end{tabular}
\end{table}


\pause 
\[\text{Average Causal Effect} =\underbrace{\E(Y^{a = 1})}_{0} - \underbrace{\E(Y^{a = 0})}_{1/2} = -1/2 \]
    

\end{frame}



\begin{frame}{Exchangeability}
$A = 1$ means vaccinated; $A = 0$ means unvaccinated\\
$Y = 1$ means covid; $Y = 0$ means no covid;
\begin{table}
\begin{tabular}[t]{c|c| c|c|c}
\toprule
& $Y^{a=1}$ & $Y^{a=0}$ & $A$ & $Y$\\
\midrule
Low Risk 1 & {\color{lightgray}0} & 0 & 0 & 0\\
Low Risk 2 &{\color{lightgray}0} & 0 & 0 & 0\\
High Risk 3 & 0 & {\color{lightgray}1} & 1 & 0\\
High Risk 4 & 0 & {\color{lightgray}1} & 1 & 0\\
\bottomrule
\end{tabular}
\end{table}

\[\text{Average Causal Effect} =\underbrace{\E(Y^{a = 1})}_{0} - \underbrace{\E(Y^{a = 0})}_{1/2} = -1/2 \]

Suppose Low Risk choose $A_i = 0$ and High Risk choose $A_i = 1$ so the potential outcomes are not independent of the observed treatment
\[\underbrace{\E(Y \mid A = 1)}_{\text{observed unvac}} - \underbrace{\E(Y \mid A = 0)}_{\text{observed unvac}} = 0 \]

\end{frame}


\begin{frame}{Exchangeability}
$A = 1$ means vaccinated; $A = 0$ means unvaccinated\\
$Y = 1$ means covid; $Y = 0$ means no covid;
\begin{table}
\begin{tabular}[t]{c|c| c|c|c}
\toprule
& $Y^{a=1}$ & $Y^{a=0}$ & $A$ & $Y$\\
\midrule
Low Risk 1 & {\color{lightgray}0} & 0 & 1 & 0\\
Low Risk 2 & 0 & {\color{lightgray}0} & 0 & 0\\
High Risk 3 & {\color{lightgray}0} & 1 & 1 & 1\\
High Risk 4 & 0 & {\color{lightgray}1} & 0 & 0\\
\bottomrule
\end{tabular}
\end{table}

\[\text{Average Causal Effect} =\underbrace{\E(Y^{a = 1})}_{0} - \underbrace{\E(Y^{a = 0})}_{1/2} = -1/2 \]
\pause 
\[\underbrace{\E(Y \mid A = 1)}_{\text{observed vac}} - \underbrace{\E(Y \mid A = 0)}_{\text{observed unvac}} = -1/2 \]






\end{frame}




\begin{frame}{Exchangeability}
In mathematical notation, exchangeability means

\vspace{1em}

\[\underbrace{Y^{a = 1}, Y^{a = 0}}_{\text{potential outcomes}} \indep \underbrace{A}_{\text{observed treatment}} \]




\end{frame}

\begin{frame}{Why is exchangeability good?}
The average causal effect (ACE) is the difference in average outcome that would occur if everyone is treated compared to the average outcome that would occur if no-one is treated   

\vspace{1em}

\[
        \text{ACE} = 
        \underbrace{\E(Y^{a = 1})}_{\text{if everyone is treated}} - \underbrace{\E(Y^{a = 0} )}_{\text{if no-one is treated}} \] 
\pause

\vspace{1em}

The problem is, we only know
\begin{itemize}
    \item  $\E(Y^{a=1} \mid A = 1)$, the average $Y^{a=1}$ among individuals who are actually treated
    \item $\E(Y^{a=0} \mid A = 0)$, the average $Y^{a=0}$ among individuals who are actually not treated
\end{itemize}


\end{frame}


\begin{frame}{Why is exchangeability good?}
When exchangeability is true, it implies

\[\underbrace{\E(Y^{a = 1}  \mid A = 1)}_{\text{Within treated}} = \underbrace{\E(Y^{a = 1}  \mid A = 0)}_{\text{Within not treated}} = \underbrace{\E(Y^{a = 1} )}_{\text{everyone}}  \]


\pause

\vspace{2em}
This allows us to identify the average causal effect (ACE)
\[
        \text{ATE} = 
        \underbrace{\E(Y^{a = 1} )}_{\text{if everyone is treated}} - \underbrace{\E(Y^{a = 0} )}_{\text{if no-one is treated}} \] because we can plug-in 
        \[\underbrace{\E(Y^{a = 1} \mid A = 1)}_{\substack{\text{outcomes for people who} \\ \text{are \textbf{actually} treated}}} \quad \text{ 
        and } \quad  \underbrace{\E(Y^{a = 0} \mid A = 0)}_{\substack{\text{outcomes for people who} \\ \text{are \textbf{actually} not treated}}}\]

\end{frame}





\goalsframe

\end{document}

